\section{周三}
\subsection{去客西馬尼}
\subsubsection{太26:30}
他們唱了詩，就出來往橄欖山去。
\subsubsection{路22:39}
耶穌出來，照常往橄欖山去，門徒也跟隨他。
\subsubsection{約18:1-2}
耶穌說了這話，就同門徒唱了詩出去，照常過了汲淪溪。往橄欖山去，門徒也跟隨他。在那裡有一個園子，他和門徒進去了。 2 賣耶穌的猶大也知道那地方，因為耶穌和門徒屢次上那裡去聚集。

\subsection{去客西馬尼路上二次預言彼得不認主 太26:30-35}
耶穌出來，照常往橄欖山去，門徒也跟隨他。那時，耶穌對他們說：今夜，你們為我的緣故都要跌倒。因為經上記著說：我要擊打牧人，羊就分散了。但我復活以後，要在你們以先往加利利去。
彼得說：眾人雖然為你的緣故跌倒，我卻永不跌倒。耶穌說：我實在告訴你，今夜雞叫以先，你要三次不認我。彼得說：我就是必須和你同死，也總不能不認你。眾門徒都是這樣說。


\subsection{耶穌的禱告-太26:36-46/可14:32-42/路22:40-46}
\subsubsection{第一次禱告}
\paragraph{太26:36-41}
耶穌同門徒來到一個地方，名叫客西馬尼，就對他們說：你們坐在這裡，等我到那邊去禱告。於是帶著彼得和西庇太的兩個兒子同去，就憂愁起來，極其難過，便對他們說：我心裡甚是憂傷，幾乎要死；你們在這裡等候，和我一同警醒。他就稍往前走，俯伏在地，禱告說：我父啊，倘若可行，求你叫這杯離開我。然而，不要照我的意思，只要照你的意思。來到門徒那裡，見他們睡著了，就對彼得說：怎麼樣？你們不能同我警醒片時嗎？
總要警醒禱告，免得入了迷惑。你們心靈固然願意，肉體卻軟弱了。
\paragraph{可14:32-38}
他們來到一個地方，名叫客西馬尼。耶穌對門徒說：你們坐在這裡，等我禱告。於是帶著彼得、雅各、約翰同去，就驚恐起來，極其難過，對他們說：我心裡甚是憂傷，幾乎要死；你們在這裡等候，警醒。他就稍往前走，俯伏在地，禱告說：倘若可行，便叫那時候過去。他說：阿爸！父啊！在你凡事都能；求你將這杯撤去。然而，不要從我的意思，只要從你的意思。耶穌回來，見他們睡著了，就對彼得說：西門，你睡覺嗎？不能警醒片時嗎？總要警醒禱告，免得入了迷惑。你們心靈固然願意，肉體卻軟弱了。
\paragraph{路22:40-46}
到了那地方，就對他們說：你們要禱告，免得入了迷惑。於是離開他們約有扔一塊石頭那麼遠，跪下禱告，說：父啊！你若願意，就把這杯撤去；然而，不要成就我的意思，只要成就你的意思。有一位天使從天上顯現，加添他的力量。耶穌極其傷痛，禱告更加懇切，汗珠如大血點滴在地上。禱告完了，就起來，到門徒那裡，見他們因為憂愁都睡著了，就對他們說：你們為什麼睡覺呢？起來禱告，免得入了迷惑！
\subsubsection{第二次禱告}
\paragraph{太26:42-44a}
第二次又去禱告說：我父啊，這杯若不能離開我，必要我喝，就願你的意旨成全。又來，見他們睡著了，因為他們的眼睛困倦。
耶穌又離開他們去了。
\paragraph{可14:39-40}
耶穌又去禱告，說的話還是與先前一樣，又來見他們睡著了，因為他們的眼睛甚是困倦；他們也不知道怎麼回答。

\subsubsection{第三次禱告}
\paragraph{太26:44b-46}
第三次禱告，說的話還是與先前一樣。於是來到門徒那裡，對他們說：現在你們仍然睡覺安歇吧（吧：或作嗎？）！時候到了，人子被賣在罪人手裡了。起來！我們走吧。看哪，賣我的人近了。

\paragraph{可14:41-42}
第三次來，對他們說：現在你們仍然睡覺安歇吧（或作嗎？）！夠了，時候到了。看哪，人子被賣在罪人手裡了。起來！我們走吧。看哪，那賣我的人近了。
\subsubsection{综合}
他們來到一個地方，名叫客西馬尼，耶穌對門徒說：「你們坐在這裡，等我禱告。」33 於是帶著彼得、雅各、約翰同去，就驚恐（憂愁）起來，極其難過， 34 對他們說：「我心裡甚是憂傷，幾乎要死。你們在這裡等候禱告，和我一同警醒，免得入了迷惑。」 35 他就稍往前走離開他們約有扔一塊石頭那麼遠，俯伏在地，禱告說倘若可行，便叫那時候過去。 36 他說：「阿爸，父啊！在你凡事都能，求你將這杯撤去！倘若可行，便叫那時候過去。然而，不要從我的意思，只要從你的意思。」有一位天使從天上顯現，加添他的力量。 44 耶穌極其傷痛，禱告更加懇切，汗珠如大血點滴在地上。37 耶穌回來到門徒那裡，見他們睡著了，就對彼得說：「西門，你睡覺嗎？你們不能同我警醒片時嗎？ 38 總要警醒禱告，免得入了迷惑。你們心靈固然願意，肉體卻軟弱了。」 39 耶穌第二次又去禱告，說的話還是與先前一樣，說：「我父啊！這杯若不能離開我，必要我喝，就願你的意旨成全！」 40 禱告完了又來，見他們因為憂愁都睡著了，因為他們的眼睛甚是困倦；他們也不知道怎麼回答。就對他們說：「你們為什麼睡覺呢？起來禱告，免得入了迷惑。」耶穌又離開他們去了。第三次禱告，說的話還是與先前一樣。 41 第三次來，對他們說：「現在你們仍然睡覺安歇吧 a！夠了，時候到了。看哪，人子被賣在罪人手裡了！ 42 起來，我們走吧！看哪，那賣我的人近了！」

\subsection{被抓}%（太26：47-56/可14：43-52/路22：47-53/約18：2-12）}
\subsubsection{太26:47-56}
說話之間，那十二個門徒裡的猶大來了，並有許多人帶著刀棒，從祭司長和民間的長老那裡與他同來。那賣耶穌的給了他們一個暗號，說：我與誰親嘴，誰就是他。你們可以拿住他。猶大隨即到耶穌跟前，說：請拉比安，就與他親嘴。耶穌對他說：朋友，你來要做的事，就做吧。於是那些人上前，下手拿住耶穌。有跟隨耶穌的一個人伸手拔出刀來，將大祭司的僕人砍了一刀，削掉了他一個耳朵。耶穌對他說：收刀入鞘吧！凡動刀的，必死在刀下。你想，我不能求我父現在為我差遣十二營多天使來嗎？若是這樣，經上所說，事情必須如此的話怎麼應驗呢？當時，耶穌對眾人說：你們帶著刀棒出來拿我，如同拿強盜嗎？我天天坐在殿裡教訓人，你們並沒有拿我。但這一切的事成就了，為要應驗先知書上的話。當下，門徒都離開他逃走了。

\subsubsection{可14:43-52}
說話之間，忽然那十二個門徒裡的猶大來了，並有許多人帶著刀棒，從祭司長和文士並長老那裡與他同來。
賣耶穌的人曾給他們一個暗號，說：我與誰親嘴，誰就是他。你們把他拿住，牢牢靠靠的帶去。
猶大來了，隨即到耶穌跟前，說：拉比，便與他親嘴。
他們就下手拿住他。
旁邊站著的人，有一個拔出刀來，將大祭司的僕人砍了一刀，削掉了他一個耳朵。
耶穌對他們說：你們帶著刀棒出來拿我，如同拿強盜嗎？
我天天教訓人，同你們在殿裡，你們並沒有拿我。但這事成就，為要應驗經上的話。
門徒都離開他，逃走了。
有一個少年人，赤身披著一塊麻布，跟隨耶穌，眾人就捉拿他。
他卻丟了麻布，赤身逃走了。


\subsubsection{路22:47-53}
說話之間，來了許多人。那十二個門徒裡名叫猶大的，走在前頭，就近耶穌，要與他親嘴。耶穌對他說：猶大！你用親嘴的暗號賣人子嗎？左右的人見光景不好，就說：主啊！我們拿刀砍可以不可以？內中有一個人把大祭司的僕人砍了一刀，削掉了他的右耳。耶穌說：到了這個地步，由他們吧！就摸那人的耳朵，把他治好了。耶穌對那些來拿他的祭司長和守殿官並長老說：你們帶著刀棒出來拿我，如同拿強盜嗎？我天天同你們在殿裡，你們不下手拿我。現在卻是你們的時候，黑暗掌權了。

\subsubsection{約18:2-12}
耶穌說了這話，就同門徒出去，過了汲淪溪。在那裡有一個園子，他和門徒進去了。賣耶穌的猶大也知道那地方，因為耶穌和門徒屢次上那裡去聚集。猶大領了一隊兵，和祭司長並法利賽人的差役，拿著燈籠、火把、兵器，就來到園裡。耶穌知道將要臨到自己的一切事，就出來對他們說：你們找誰？他們回答說：找拿撒勒人耶穌。耶穌說：我就是。賣他的猶大也同他們站在那裡。耶穌一說我就是，他們就退後倒在地上。他又問他們說：你們找誰？他們說：找拿撒勒人耶穌。耶穌說：我已經告訴你們，我就是。你們若找我，就讓這些人去吧。這要應驗耶穌從前的話，說：你所賜給我的人，我沒有失落一個。西門彼得帶著一把刀，就拔出來，將大祭司的僕人砍了一刀，削掉他的右耳；那僕人名叫馬勒古。耶穌就對彼得說：收刀入鞘吧，我父所給我的那杯，我豈可不喝呢？那隊兵和千夫長，並猶太人的差役就拿住耶穌，把他捆綁了，
\subsubsection{综合}
說話之間，忽然那十二個門徒裡的猶大來了走在前頭，並有許多人帶著刀棒，從祭司長和文士並長老那裡與他同來。 44 賣耶穌的人曾給他們一個暗號，說：「我與誰親嘴，誰就是他。你們把他拿住，牢牢靠靠地帶去。」 45 猶大來了，隨即到耶穌跟前，說：「請拉比安！」便與他親嘴。耶穌對他說：「猶大，你用親嘴的暗號賣人子嗎？朋友，你來要做的事就做吧！」 46 他們就下手拿住他。左右的人見光景不好，就說：「主啊，我們拿刀砍可以不可以？」 47 旁邊站著的人，有一個跟隨耶穌的拔出刀來，將大祭司的僕人砍了一刀，削掉了他右耳朵。耶穌對他說：「收刀入鞘吧！凡動刀的，必死在刀下。 53 你想我不能求我父現在為我差遣十二營多天使來嗎？ 54 若是這樣，經上所說事情必須如此的話怎麼應驗呢？到了這個地步，由他們吧！」就摸那人的耳朵，把他治好了。 55 當時，耶穌對那些來拿他的祭司長和守殿官並長老說：「你們帶著刀棒出來拿我，如同拿強盜嗎？ 49 我天天教訓人，同你們在殿裡，你們並沒有拿我。但這事成就，為要應驗經上的話。現在卻是你們的時候，黑暗掌權了。」 50 門徒都離開他逃走了。有一個少年人赤身披著一塊麻布，跟隨耶穌，眾人就捉拿他。 52 他卻丟了麻布，赤身逃走了。


\subsection{被猶太人審問}
\subsubsection{在亞拿府被審問-路22：54/約18：13-14，19-23}

\begin{table}[htbp]
  \centering 
    \begin{tabular}{|p{3em}|p{39em}|}
    \toprule
    \multicolumn{2}{|p{42em}|}{\textcolor[rgb]{ 0,  .075,  .125}{那隊兵和千夫長並猶太人的差役就拿住耶穌，把他捆綁了，先帶到亞那面前（彼得初次否認主）}} \\
    \midrule
    \textcolor[rgb]{ 0,  .075,  .125}{大祭司} & \textcolor[rgb]{ 0,  .075,  .125}{大祭司就以耶穌的門徒和他的教訓盤問他。} \\
    \midrule
    \textcolor[rgb]{ 0,  .075,  .125}{耶穌} & 我從來是明明地對世人說話。我常在會堂和殿裡，就是猶太人聚集的地方教訓人，我在暗地裡並沒有說什麼。你為什麼問我呢？可以問那聽見的人我對他們說的是什麼，我所說的，他們都知道。 \\
    \midrule
    \textcolor[rgb]{ 0,  .075,  .125}{差役} & 旁邊站著的一個差役用手掌打他說：你這樣回答大祭司嗎 \\
    \midrule
    \textcolor[rgb]{ 0,  .075,  .125}{耶穌} & 我若說的不是，你可以指證那不是。我若說得是，你為什麼打我呢？ \\
    \bottomrule
    \end{tabular}%
     \caption{在亞那前受審 （約18:13-14，19-23）}
  \label{tab:addlabel}%
\end{table}


\paragraph{路22:54}
他們拿住耶穌，把他帶到大祭司的宅裡。彼得遠遠的跟著。

\paragraph{約18:13-14，19-23}
先帶到亞那面前，因為亞那是本年作大祭司該亞法的岳父。這該亞法就是從前向猶太人發議論說一個人替百姓死是有益的那位。大祭司就以耶穌的門徒和他的教訓盤問他。耶穌回答說：我從來是明明的對世人說話。我常在會堂和殿裡，就是猶太人聚集的地方教訓人；我在暗地裡並沒有說什麼。你為什麼問我呢？可以問那聽見的人，我對他們說的是什麼；我所說的，他們都知道。耶穌說了這話，旁邊站著的一個差役用手掌打他，說：你這樣回答大祭司嗎？耶穌說：我若說的不是，你可以指證那不是；我若說的是，你為什麼打我呢？

\subsubsection{彼得第一次不認主-太26:69-70/可14:66-68/路22:55-57/約18:15-18}
\paragraph{太26:69-70}彼得在外面院子裡坐著，有一個使女前來，說：你素來也是同那加利利人耶穌一夥的。彼得在眾人面前卻不承認，說：我不知道你說的是什麼！

\paragraph{可14:66-68}
彼得在下邊院子裡；來了大祭司的一個使女，
見彼得烤火，就看著他，說：你素來也是同拿撒勒人耶穌一夥的。
彼得卻不承認，說：我不知道，也不明白你說的是什麼。於是出來，到了前院，雞就叫了。

\paragraph{路22:55-57}
他們在院子裡生了火，一同坐著；彼得也坐在他們中間。有一個使女看見彼得坐在火光裡，就定睛看他，說：這個人素來也是同那人一夥的。彼得卻不承認，說：女子，我不認得他。


\paragraph{約18:15-18}
西門彼得跟著耶穌，還有一個門徒跟著。那門徒是大祭司所認識的，他就同耶穌進了大祭司的院子。彼得卻站在門外。大祭司所認識的那個門徒出來，和看門的使女說了一聲，就領彼得進去。那看門的使女對彼得說：你不也是這人的門徒嗎？他說：我不是。僕人和差役因為天冷，就生了炭火，站在那裡烤火；彼得也同他們站著烤火。




\subsubsection{在該亞法府被審問-太26:57-68/可14:53-65/路22:63-71/約18:24}
\begin{table}[htbp]
  \centering 
    \begin{tabular}{|p{3em}|p{40em}|}
    \toprule
    \multicolumn{2}{|p{43em}|}{\textcolor[rgb]{ 0,  .075,  .125}{亞那就把耶穌解到大祭司該亞法那裡，是捆著解去的。文士和長老已經在那裡聚會。祭司長和全公會尋找假見證控告耶穌，要治死他。雖有好些人來作假見證，總得不著實據。他們的見證各不相合}} \\
    \midrule
    \textcolor[rgb]{ 0,  .075,  .125}{兩個人} & \textcolor[rgb]{ 0,  .075,  .125}{末後，有兩個人前來說：「這個人曾說：『我能拆毀神的殿，三日內又建造起來。』」他們就是這麼作見證，也是各不相合。(太26:60b-61,可14:57-58)} \\
    \midrule
    \textcolor[rgb]{ 0,  .075,  .125}{大祭司} & \textcolor[rgb]{ 0,  .075,  .125}{大祭司就站起來，對耶穌說：「你什麼都不回答嗎？這些人作見證告你的是什麼呢？」(太26:62, 可14:60)} \\
    \midrule
    \textcolor[rgb]{ 0,  .075,  .125}{耶穌} & \textcolor[rgb]{ 0,  .075,  .125}{耶穌卻不言語。一句也不回答(太26:63a,可14:61a)} \\
    \midrule
    \textcolor[rgb]{ 0,  .075,  .125}{大祭司} & \textcolor[rgb]{ 0,  .075,  .125}{我指著永生神叫你起誓告訴我們，你是神的兒子基督不是？(太26:63b, 可14:61b)} \\
    \midrule
    \textcolor[rgb]{ 0,  .075,  .125}{耶穌} & 你說的是。然而我告訴你們：後來你們要看見人子坐在那權能者的右邊，駕著天上的雲降臨。(太26:64,可14:62) \\
    \midrule
    \textcolor[rgb]{ 0,  .075,  .125}{大祭司} & 撕開衣服說：「他說了僭妄的話！我們何必再用見證人呢？這僭妄的話現在你們都聽見了。你們的意見如何？」 \\
    \midrule
    \textcolor[rgb]{ 0,  .075,  .125}{他們} & 回答說：「他是該死的。」就吐唾沫在他臉上，又蒙著他的臉，用拳頭打他，也有用手掌打他的，說：「基督啊！你是先知，告訴我們打你的是誰？」「你說預言吧！」 \\
    \midrule
    差役&	差役接過他來打他。看守的人戲弄他，還用許多別的話辱罵他。\\
    \bottomrule
    \end{tabular}%
     \caption{在大祭司該亞法前受審}
  \label{tab:addlabel}%
\end{table}
\paragraph{太26:57-68} 
拿耶穌的人把他帶到大祭司該亞法那裡去；文士和長老已經在那裡聚會。彼得遠遠的跟著耶穌，直到大祭司的院子，進到裡面，就和差役同坐，要看這事到底怎樣。祭司長和全公會尋找假見證控告耶穌，要治死他。雖有好些人來作假見證，總得不著實據。末後有兩個人前來，說：這個人曾說：我能拆毀神的殿，三日內又建造起來。大祭司就站起來，對耶穌說：你什麼都不回答嗎？這些人作見證告你的是什麼呢？耶穌卻不言語。大祭司對他說：我指著永生神叫你起誓告訴我們，你是神的兒子基督不是？耶穌對他說：你說的是。然而，我告訴你們，後來你們要看見人子坐在那權能者的右邊，駕著天上的雲降臨。大祭司就撕開衣服，說：他說了僭妄的話，我們何必再用見證人呢？這僭妄的話，現在你們都聽見了。你們的意見如何？他們回答說：他是該死的。他們就吐唾沫在他臉上，用拳頭打他；也有用手掌打他的，說：基督啊！你是先知，告訴我們打你的是誰？


\paragraph{可14:53-65}
他們把耶穌帶到大祭司那裡，又有眾祭司長和長老並文士都來和大祭司一同聚集。
彼得遠遠的跟著耶穌，一直進入大祭司的院裡，和差役一同坐在火光裡烤火。
祭司長和全公會尋找見證控告耶穌，要治死他，卻尋不著。
因為有好些人作假見證告他，只是他們的見證各不相合。
又有幾個人站起來作假見證告他，說：
我們聽見他說：我要拆毀這人手所造的殿，三日內就另造一座不是人手所造的。
他們就是這麼作見證，也是各不相合。
大祭司起來站在中間，問耶穌說：你什麼都不回答嗎？這些人作見證告你的是什麼呢？
耶穌卻不言語，一句也不回答。大祭司又問他說：你是那當稱頌者的兒子基督不是？
耶穌說：我是。你們必看見人子坐在那權能者的右邊，駕著天上的雲降臨。
大祭司就撕開衣服，說：我們何必再用見證人呢？
你們已經聽見他這僭妄的話了。你們的意見如何？他們都定他該死的罪。
就有人吐唾沫在他臉上，又蒙著他的臉，用拳頭打他，對他說：你說預言吧！差役接過他來，用手掌打他。

\paragraph{路22:63-65}
看守耶穌的人戲弄他，打他，又蒙著他的眼，問他說：你是先知，告訴我們打你的是誰？他們還用許多別的話辱罵他。



\paragraph{約18:24}
亞那就把耶穌解到大祭司該亞法那裡，仍是捆著解去的。


\subsubsection{彼得二，三次不認主-太26:71-75/可14:69-72/路22:58-62/約18:25-27}
\paragraph{太26:71-75}
既出去，到了門口，又有一個使女看見他，就對那裡的人說：這個人也是同拿撒勒人耶穌一夥的。彼得又不承認，並且起誓說：我不認得那個人。過了不多的時候，旁邊站著的人前來，對彼得說：你真是他們一黨的，你的口音把你露出來了。
彼得就發咒起誓的說：我不認得那個人。立時，雞就叫了。彼得想起耶穌所說的話：雞叫以先，你要三次不認我。他就出去痛哭。

\paragraph{可14:69-72}
那使女看見他，又對旁邊站著的人說：這也是他們一黨的。彼得又不承認。過了不多的時候，旁邊站著的人又對彼得說：你真是他們一黨的！因為你是加利利人。彼得就發咒起誓的說：我不認得你們說的這個人。立時雞叫了第二遍。彼得想起耶穌對他所說的話：雞叫兩遍以先，你要三次不認我。思想起來，就哭了。

\paragraph{路22:58-62}
過了不多的時候，又有一個人看見他，說：你也是他們一黨的。彼得說：你這個人！我不是。約過了一小時，又有一個人極力的說：他實在是同那人一夥的，因為他也是加利利人。彼得說：你這個人！我不曉得你說的是什麼！正說話之間，雞就叫了。主轉過身來看彼得。彼得便想起主對他所說的話：今日雞叫以先，你要三次不認我。他就出去痛哭。

\paragraph{約18:25-27}
西門彼得正站著烤火，有人對他說：你不也是他的門徒嗎？彼得不承認，說：我不是。有大祭司的一個僕人，是彼得削掉耳朵那人的親屬，說：我不是看見你同他在園子裡嗎？彼得又不承認。立時雞就叫了。

\subsubsection{彼得三次不認主综述}
西門彼得遠遠地跟著耶穌，還有一個門徒跟著。那門徒是大祭司所認識的，他就同耶穌進了大祭司的院子， 16 彼得卻站在門外。大祭司所認識的那個門徒出來，和看門的使女說了一聲，就領彼得進去,  彼得在外面院子裡就和差役同坐，要看這事到底怎樣。僕人和差役因為天冷，就生了炭火，站在那裡烤火，彼得也同他們站著烤火。 17 那看門的使女見彼得烤火，就定睛看著他對彼得說：「你不也是這人的門徒嗎？你素來也是同那加利利人耶穌一夥的。」彼得卻不承認，說：「女子，我不是。我不知道你說的是什麼。」 18 既出去，到了門口，過了不多的時候，又有一個使女看見他，就對那裡站著的的人說：「這個人也是同拿撒勒人耶穌一夥的。」 72 彼得又不承認，並且起誓說：「你這個人！我不認得那個人！」 73 過了不多的時候, 約過了一小時，旁邊站著的人前來，極力地對彼得說：「你真是他們一黨的，你的口音把你露出來了。因為你是加利利人。」 74 彼得就發咒起誓地說：「你這個人！我不曉得你說的是什麼！我不認得那個人！」立時正說話之間，雞就叫了第二遍。 75 主轉過身來看彼得，彼得想起耶穌所說的話：「雞叫兩遍以先，你要三次不認我。」思想起來，他就出去痛哭。

\subsubsection{在亞那前受審}
那隊兵和千夫長並猶太人的差役就拿住耶穌，把他捆綁了， 13 先帶到亞那面前，因為亞那是本年做大祭司該亞法的岳父。 14 這該亞法就是從前向猶太人發議論說「一個人替百姓死是有益的」那位。大祭司就以耶穌的門徒和他的教訓盤問他。 20 耶穌回答說：「我從來是明明地對世人說話。我常在會堂和殿裡，就是猶太人聚集的地方教訓人，我在暗地裡並沒有說什麼。21 你為什麼問我呢？可以問那聽見的人我對他們說的是什麼，我所說的，他們都知道。」22 耶穌說了這話，旁邊站著的一個差役用手掌打他，說：「你這樣回答大祭司嗎？」 23 耶穌說：「我若說的不是，你可以指證那不是。我若說得是，你為什麼打我呢？」24 亞那就把耶穌解到大祭司該亞法那裡，仍是捆著解去的。

\subsubsection{在大祭司該亞法前受審}
他們把耶穌帶到大祭司該亞法那裡，又有眾祭司長和長老並文士都來和大祭司一同聚集。  55 祭司長和全公會尋找見證控告耶穌，要治死他，卻尋不著。 56 因為有好些人作假見證告他，只是他們的見證各不相合總得不著實據。末後，又有幾個人站起來作假見證告他，說： 58 「我們聽見他說：『我要拆毀這人手所造的神的殿，三日內就另造一座不是人手所造的。』」 59 他們就是這麼作見證，也是各不相合。 60 大祭司起來站在中間，問耶穌說：「你什麼都不回答嗎？這些人作見證告你的是什麼呢？」 61 耶穌卻不言語，一句也不回答。大祭司又問他說：「我指著永生神叫你起誓告訴我們，你是那當稱頌者的兒子基督不是？」 62 耶穌說：「我是。你說的是。然而我告訴你們：你們必看見人子坐在那權能者的右邊，駕著天上的雲降臨。」 63 大祭司就撕開衣服，說：「我們何必再用見證人呢？ 64 你們已經聽見他這僭妄的話了。你們的意見如何？」他們回答說：「他是該死的。」他們都定他該死的罪。 65 看守耶穌的人戲弄他，打他，就有人吐唾沫在他臉上，又蒙著他的臉，用拳頭打他，對他說：「你說預言吧！基督啊！你是先知，告訴我們打你的是誰？」差役接過他來，用手掌打他。他們還用許多別的話辱罵他。

\subsubsection{在公會被審問-太27:1-2/可15:1/路22:66-71}

\begin{table}[htbp]
  \centering
  
    \begin{tabular}{|p{2em}|p{42em}|}
    \toprule
    \multicolumn{2}{|p{44em}|}{一到早晨，祭司長和長老、文士、全公會的人大家商議，就把耶穌捆綁，解去交給彼拉多。} \\
    \midrule
    \textcolor[rgb]{ 0,  .075,  .125}{眾人} & 「你若是基督，就告訴我們。」 \\
    \midrule
    \textcolor[rgb]{ 0,  .075,  .125}{耶穌} & 我若告訴你們，你們也不信；我若問你們，你們也不回答。從今以後，人子要坐在神權能的右邊。 \\
    \midrule
    \textcolor[rgb]{ 0,  .075,  .125}{他們} & 這樣，你是神的兒子嗎？ \\
    \midrule
    \textcolor[rgb]{ 0,  .075,  .125}{耶穌} & 你們所說的是。 \\
    \midrule
    \textcolor[rgb]{ 0,  .075,  .125}{他們} & \textcolor[rgb]{ 0,  .075,  .125}{何必再用見證呢？他親口所說的，我們都親自聽見了。} \\
    \bottomrule
    \end{tabular}%
    \caption{在公會前受審}
  \label{tab:addlabel}%
\end{table}

\paragraph{太27:1-2}
到了早晨，眾祭司長和民間的長老大家商議要治死耶穌，就把他捆綁，解去，交給巡撫彼拉多。

\paragraph{可15:1}
一到早晨，祭司長和長老、文士、全公會的人大家商議，就把耶穌捆綁，解去交給彼拉多。

\paragraph{路22:66-71}
天一亮，民間的眾長老連祭司長帶文士都聚會，把耶穌帶到他們的公會裡，說：你若是基督，就告訴我們。耶穌說：我若告訴你們，你們也不信；我若問你們，你們也不回答。從今以後，人子要坐在神權能的右邊。他們都說：這樣，你是神的兒子嗎？耶穌說：你們所說的是。他們說：何必再用見證呢？他親口所說的，我們都親自聽見了。
\paragraph{综述}
天一亮，民間的眾長老連祭司長帶文士都聚會，把耶穌帶到他們的公會裡， 67 說：「你若是基督，就告訴我們。」耶穌說：「我若告訴你們，你們也不信； 68 我若問你們，你們也不回答。 69 從今以後，人子要坐在神權能的右邊。」 70 他們都說：「這樣，你是神的兒子嗎？」耶穌說：「你們所說的是。」 71 他們說：「何必再用見證呢？他親口所說的，我們都親自聽見了。」

\subsubsection{猶大自縊-太27:3-10}
這時候，賣耶穌的猶大看見耶穌已經定了罪，就後悔，把那三十塊錢拿回來給祭司長和長老，說：我賣了無辜之人的血是有罪了。他們說：那與我們有什麼相干？你自己承當吧！猶大就把那銀錢丟在殿裡，出去吊死了。祭司長拾起銀錢來，說：這是血價，不可放在庫裡。他們商議，就用那銀錢買了窰戶的一塊田，為要埋葬外鄉人。所以那塊田直到今日還叫做血田。這就應驗了先知耶利米的話，說：他們用那三十塊錢，就是被估定之人的價錢，是以色列人中所估定的，買了窰戶的一塊田；這是照著主所吩咐我的。


\subsection{被外邦人審問}
\subsubsection{被巡撫彼拉多審問-太27:11-14/可15:1-5/路23:1-6/約18:28-38a}
\paragraph{太27:11-14}耶穌站在巡撫面前；巡撫問他說：你是猶太人的王嗎？耶穌說：你說的是。
他被祭司長和長老控告的時候，什麼都不回答。
彼拉多就對他說：他們作見證告你這麼多的事，你沒有聽見嗎？
耶穌仍不回答，連一句話也不說，以致巡撫甚覺希奇。



\paragraph{可15:2-5}
彼拉多問他說：你是猶太人的王嗎？耶穌回答說：你說的是。
祭司長告他許多的事。
彼拉多又問他說：你看，他們告你這麼多的事，你什麼都不回答嗎？
耶穌仍不回答，以致彼拉多覺得希奇。




\paragraph{路23:1-6}眾人都起來，把耶穌解到彼拉多面前，就告他說：我們見這人誘惑國民，禁止納稅給該撒，並說自己是基督，是王。彼拉多問耶穌說：你是猶太人的王嗎？耶穌回答說：你說的是。彼拉多對祭司長和眾人說：我查不出這人有什麼罪來。但他們越發極力的說：他煽惑百姓，在猶太遍地傳道，從加利利起，直到這裡了。彼拉多一聽見，就問：這人是加利利人嗎？

\paragraph{約18：28-38a}
眾人將耶穌從該亞法那裡往衙門內解去，那時天還早。他們自己卻不進衙門，恐怕染了污穢，不能吃逾越節的筵席。彼拉多就出來，到他們那裡，說：你們告這人是為什麼事呢？
他們回答說：這人若不是作惡的，我們就不把他交給你。彼拉多說：你們自己帶他去，按著你們的律法審問他吧。猶太人說：我們沒有殺人的權柄。這要應驗耶穌所說自己將要怎樣死的話了。
彼拉多又進了衙門，叫耶穌來，對他說：你是猶太人的王嗎？
耶穌回答說：這話是你自己說的，還是別人論我對你說的呢？
彼拉多說：我豈是猶太人呢？你本國的人和祭司長把你交給我。你做了什麼事呢？耶穌回答說：我的國不屬這世界；我的國若屬這世界，我的臣僕必要爭戰，使我不至於被交給猶太人。只是我的國不屬這世界。彼拉多就對他說：這樣，你是王嗎？耶穌回答說：你說我是王。我為此而生，也為此來到世間，特為給真理作見證。凡屬真理的人就聽我的話。彼拉多說：真理是什麼呢？
\paragraph{综述}
一到早晨，祭司長和長老、文士、全公會的人大家商議要治死耶穌，就把耶穌捆綁，眾人將耶穌從該亞法那裡往衙門內解去，交給彼拉多。那時天還早。他們自己卻不進衙門，恐怕染了汙穢，不能吃逾越節的筵席。 29 彼拉多就出來，到他們那裡，說：「你們告這人是為什麼事呢？」 30 他們回答說：「這人若不是作惡的，我們就不把他交給你。我們見這人誘惑國民，禁止納稅給愷撒，並說自己是基督，是王。」 31 彼拉多說：「你們自己帶他去，按著你們的律法審問他吧！」猶太人說：「我們沒有殺人的權柄。」 32 這要應驗耶穌所說自己將要怎樣死的話了。
 2 彼拉多又進了衙門，叫耶穌來站在巡撫面前，巡撫問他說：「你是猶太人的王嗎？」耶穌回答說：「這話是你自己說的，還是別人論我對你說的呢？」 35 彼拉多說：「我豈是猶太人呢？你本國的人和祭司長把你交給我。你做了什麼事呢？」 36 耶穌回答說：「我的國不屬這世界。我的國若屬這世界，我的臣僕必要爭戰，使我不至於被交給猶太人。只是我的國不屬這世界。」 37 彼拉多就對他說：「這樣，你是王嗎？」耶穌回答說：「你說我是王。你說的是。我為此而生，也為此來到世間，特為給真理作見證。凡屬真理的人就聽我的話。」彼拉多說：「真理是什麼呢？」3 祭司長和長老控告他許多的事。耶穌什麼都不回答。 4 彼拉多又問他說：「你看，他們告你這麼多的事，你沒有聽見嗎？你什麼都不回答嗎？」 5 耶穌仍不回答，連一句話也不說，以致彼拉多覺得稀奇。
彼拉多又出來到猶太人那裡，對他們說：「我查不出他有什麼罪來。 39 但你們有個規矩，在逾越節要我給你們釋放一個人，你們要我給你們釋放猶太人的王嗎？」但他們越發極力地說：「他煽惑百姓，在猶太遍地傳道，從加利利起直到這裡了。」 6 彼拉多一聽見，就問：「這人是加利利人嗎？」 

\subsubsection{被希律審問-路23：7-12}
既曉得耶穌屬希律所管，就把他送到希律那裡去。那時希律正在耶路撒冷。希律看見耶穌，就很歡喜；因為聽見過他的事，久已想要見他，並且指望看他行一件神蹟。於是問他許多的話；耶穌卻一言不答。祭司長和文士都站著，極力的告他。希律和他的兵丁就藐視耶穌，戲弄他，給他穿上華麗衣服，把他送回彼拉多那裡去。從前希律和彼拉多彼此有仇，在那一天就成了朋友。


\subsubsection{再次被巡撫彼拉多審問-太27：15-31/可15：6-20/路23：13-25/約18：38b-19:12}
\paragraph{太27：15-31}巡撫有一個常例，每逢這節期，隨眾人所要的釋放一個囚犯給他們。
當時有一個出名的囚犯叫巴拉巴。
眾人聚集的時候，彼拉多就對他們說：你們要我釋放哪一個給你們？是巴拉巴呢？是稱為基督的耶穌呢？
巡撫原知道他們是因為嫉妒才把他解了來。
正坐堂的時候，他的夫人打發人來說：這義人的事，你一點不可管，因為我今天在夢中為他受了許多的苦。
祭司長和長老挑唆眾人，求釋放巴拉巴，除滅耶穌。
巡撫對眾人說：這兩個人，你們要我釋放哪一個給你們呢？他們說：巴拉巴。
彼拉多說：這樣，那稱為基督的耶穌我怎麼辦他呢？他們都說：把他釘十字架！
巡撫說：為什麼呢？他做了什麼惡事呢？他們便極力的喊著說：把他釘十字架！
彼拉多見說也無濟於事，反要生亂，就拿水在眾人面前洗手，說：流這義人的血，罪不在我，你們承當吧。
眾人都回答說：他的血歸到我們和我們的子孫身上。
於是彼拉多釋放巴拉巴給他們，把耶穌鞭打了，交給人釘十字架。
巡撫的兵就把耶穌帶進衙門，叫全營的兵都聚集在他那裡。
他們給他脫了衣服，穿上一件朱紅色袍子，
用荊棘編做冠冕，戴在他頭上，拿一根葦子放在他右手裡，跪在他面前，戲弄他，說：恭喜，猶太人的王啊！
又吐唾沫在他臉上，拿葦子打他的頭。
戲弄完了，就給他脫了袍子，仍穿上他自己的衣服，帶他出去，要釘十字架。

\paragraph{可15：6-20}每逢這節期，巡撫照眾人所求的，釋放一個囚犯給他們。
有一個人名叫巴拉巴，和作亂的人一同捆綁。他們作亂的時候，曾殺過人。
眾人上去求巡撫，照常例給他們辦。
彼拉多說：你們要我釋放猶太人的王給你們嗎？
他原曉得，祭司長是因為嫉妒才把耶穌解了來。
只是祭司長挑唆眾人，寧可釋放巴拉巴給他們。
彼拉多又說：那麼樣，你們所稱為猶太人的王，我怎麼辦他呢？
他們又喊著說：把他釘十字架！
彼拉多說：為什麼呢？他做了什麼惡事呢？他們便極力的喊著說：把他釘十字架！
彼拉多要叫眾人喜悅，就釋放巴拉巴給他們，將耶穌鞭打了，交給人釘十字架。
兵丁把耶穌帶進衙門院裡，叫齊了全營的兵。
他們給他穿上紫袍，又用荊棘編做冠冕給他戴上，
就慶賀他說：恭喜，猶太人的王啊！
又拿一根葦子打他的頭，吐唾沫在他臉上，屈膝拜他。
戲弄完了，就給他脫了紫袍，仍穿上他自己的衣服，帶他出去，要釘十字架。

\paragraph{路23：13-25}
\begin{itemize} 
\item one \dots{彼拉多傳齊了祭司長和官府並百姓，就對他們說：你們解這人到我這裡，說他是誘惑百姓的。看哪，我也曾將你們告他的事，在你們面前審問他，並沒有查出他什麼罪來；就是希律也是如此，所以把他送回來。可見他沒有做什麼該死的事。故此，我要責打他，把他釋放了。（有古卷在此有：每逢這節期，巡撫必須釋放一個囚犯給他們。）眾人卻一齊喊著說：除掉這個人！釋放巴拉巴給我們！這巴拉巴是因在城裡作亂殺人，下在監裡的。}
\item two \dots{彼拉多願意釋放耶穌，就又勸解他們。無奈他們喊著說：釘他十字架！釘他十字架！}
\item three \dots{彼拉多第三次對他們說：為什麼呢？這人做了什麼惡事呢？我並沒有查出他什麼該死的罪來。所以，我要責打他，把他釋放了。他們大聲催逼彼拉多，求他把耶穌釘在十字架上。他們的聲音就得了勝。彼拉多這才照他們所求的定案，把他們所求的那作亂殺人、下在監裡的釋放了，把耶穌交給他們，任憑他們的意思行。}
\end{itemize}

\paragraph{約18：38b-19:12}
說了這話，又出來到猶太人那裡，對他們說：我查不出他有什麼罪來。但你們有個規矩，在逾越節要我給你們釋放一個人，你們要我給你們釋放猶太人的王嗎？他們又喊著說：不要這人，要巴拉巴！這巴拉巴是個強盜。

當下彼拉多將耶穌鞭打了。兵丁用荊棘編做冠冕戴在他頭上，給他穿上紫袍，又挨近他，說：恭喜，猶太人的王啊！他們就用手掌打他。彼拉多又出來對眾人說：我帶他出來見你們，叫你們知道我查不出他有什麼罪來。耶穌出來，戴著荊棘冠冕，穿著紫袍。彼拉多對他們說：你們看這個人！祭司長和差役看見他，就喊著說：釘他十字架！釘他十字架！彼拉多說：你們自己把他釘十字架吧！我查不出他有什麼罪來。猶太人回答說：我們有律法，按那律法，他是該死的，因他以自己為神的兒子。彼拉多聽見這話，越發害怕，
又進衙門，對耶穌說：你是哪裡來的？耶穌卻不回答。彼拉多說：你不對我說話嗎？你豈不知我有權柄釋放你，也有權柄把你釘十字架嗎？耶穌回答說：若不是從上頭賜給你的，你就毫無權柄辦我。所以，把我交給你的那人罪更重了。從此，彼拉多想要釋放耶穌，無奈猶太人喊著說：你若釋放這個人，就不是該撒的忠臣（原文作朋友）。凡以自己為王的，就是背叛該撒了。

\subsubsection{耶穌被衆人棄絕}
彼拉多傳齊了祭司長和官府並百姓， 14 就對他們說：「你們解這人到我這裡，說他是誘惑百姓的。看哪，我也曾將你們告他的事在你們面前審問他，並沒有查出他什麼罪來。 15 就是希律也是如此，所以把他送回來。可見他沒有做什麼該死的事。 16 故此，我要責打他，把他釋放了。」眾人卻一齊喊著說：「除掉這個人！釋放巴拉巴給我們！」巡撫有一個常例，每逢這節期，隨眾人所要的釋放一個囚犯給他們。 16 當時有一個出名的囚犯叫巴拉巴。和作亂的人一同捆綁；他們作亂的時候，曾殺過人。 17 眾人聚集的時候，上去求巡撫照常例給他們辦。彼拉多就對他們說：「你們要我釋放哪一個給你們？是巴拉巴呢，是稱為基督的耶穌呢？」 18 巡撫原知道他們是因為嫉妒才把他解了來。 19 正坐堂的時候，他的夫人打發人來說：「這義人的事你一點不可管，因為我今天在夢中為他受了許多的苦。」 20 只是祭司長和長老挑唆眾人，寧可求釋放巴拉巴，除滅耶穌。 21 彼拉多願意釋放耶穌，就又勸解他們。巡撫對眾人說：「這兩個人，你們要我釋放哪一個給你們呢？」他們說：「巴拉巴！」 22 彼拉多說：「這樣，那稱為基督的耶穌猶太人的王，我怎麼辦他呢？」他們都說：「把他釘十字架！」 23 巡撫說：「為什麼呢？他做了什麼惡事呢？」他們大聲催逼彼拉多，便極力地喊著說：「把他釘十字架！」 24 彼拉多見說也無濟於事，反要生亂，要叫眾人喜悅，就拿水在眾人面前洗手，說：「流這義人的血，罪不在我，你們承當吧！」 25 眾人都回答說：「他的血歸到我們和我們的子孫身上！」 26 於是彼拉多釋放巴拉巴給他們

\subsubsection{耶穌再被衆人棄絕}
彼拉多又出來對眾人說：「我帶他出來見你們，叫你們知道我查不出他有什麼罪來。」 5 耶穌出來，戴著荊棘冠冕，穿著紫袍。彼拉多對他們說：「你們看，這個人！」 6 祭司長和差役看見他，就喊著說：「釘他十字架！釘他十字架！」彼拉多說：「你們自己把他釘十字架吧，我查不出他有什麼罪來。」 7 猶太人回答說：「我們有律法，按那律法，他是該死的！因他以自己為神的兒子。」 8 彼拉多聽見這話，越發害怕， 9 又進衙門，對耶穌說：「你是哪裡來的？」耶穌卻不回答。 10 彼拉多說：「你不對我說話嗎？你豈不知我有權柄釋放你，也有權柄把你釘十字架嗎？」 11 耶穌回答說：「若不是從上頭賜給你的，你就毫無權柄辦我。所以，把我交給你的那人罪更重了。」 12 從此彼拉多想要釋放耶穌，無奈猶太人喊著說：「你若釋放這個人，就不是愷撒的忠臣 a。凡以自己為王的，就是背叛愷撒了。」 彼拉多第三次對他們說：「為什麼呢？這人做了什麼惡事呢？我並沒有查出他什麼該死的罪來。所以，我要責打他，把他釋放了。」 23 他們大聲催逼彼拉多，求他把耶穌釘在十字架上，他們的聲音就得了勝。

\subsubsection{耶穌被兵丁羞辱}
巡撫的兵就把耶穌帶進衙門，叫全營的兵都聚集在他那裡。 28 他們給他脫了衣服，穿上一件朱紅色袍子， 29 用荊棘編做冠冕，戴在他頭上，拿一根葦子放在他右手裡，跪在他面前，戲弄他，就慶賀他說：「恭喜，猶太人的王啊！」 30 又吐唾沫在他臉上，拿葦子打他的頭。屈膝拜他。 31 戲弄完了，就給他脫了袍子，仍穿上他自己的衣服，帶他出去，要釘十字架。

\subsubsection{巡撫彼拉多在華石處宣佈釘死耶穌-約19：13-16}
彼拉多聽見這話，就帶耶穌出來，到了一個地方，名叫鋪華石處，希伯來話叫厄巴大，就在那裡坐堂。
那日是預備逾越節的日子，約有午正。彼拉多對猶太人說：看哪，這是你們的王！
他們喊著說：除掉他！除掉他！釘他在十字架上！彼拉多說：我可以把你們的王釘十字架嗎？祭司長回答說：除了該撒，我們沒有王。彼拉多這才照他們所求的定案
於是彼拉多將耶穌交給他們去釘十字架。



\subsection{耶穌被交給人釘死}

\begin{table}[htbp]
  \centering
 
    \begin{tabular}{|p{13.5em}|p{4.055em}|r|r|p{4.055em}|r|}
    \toprule
    事 件   & 地 點   & \multicolumn{1}{p{4.055em}|}{馬太福音} & \multicolumn{1}{p{4.055em}|}{馬可福音} & 路加福音  & \multicolumn{1}{p{4.055em}|}{約翰福音} \\
    \midrule
    耶穌被交彼拉多審問 & 耶路撒冷  & \multicolumn{1}{p{4.055em}|}{27:1-2，27:11-14} & \multicolumn{1}{p{4.055em}|}{15:1-5} & 23:1-7 & \multicolumn{1}{p{4.055em}|}{18:28-38} \\
    \midrule
    猶大自盡  & \multicolumn{1}{r|}{} & \multicolumn{1}{p{4.055em}|}{27:3-10} &       & \multicolumn{1}{r|}{} &  \\
    \midrule
    耶穌被希律審問 & 耶路撒冷  &       &       & 23:8-12 &  \\
    \midrule
    耶穌被衆人棄絕 & 耶路撒冷  & \multicolumn{1}{p{4.055em}|}{27:15-26上} & \multicolumn{1}{p{4.055em}|}{15:6-15上} & 23:13-21 & \multicolumn{1}{p{4.055em}|}{18:39-40} \\
    \midrule
    耶穌被兵丁羞辱 & 耶路撒冷  & \multicolumn{1}{p{4.055em}|}{27:27-31} & \multicolumn{1}{p{4.055em}|}{15:16-20} & \multicolumn{1}{r|}{} & \multicolumn{1}{p{4.055em}|}{19:1-3} \\
    \midrule
    耶穌再被衆人棄絕 & 耶路撒冷  &       &       & 23:22-23 & \multicolumn{1}{p{4.055em}|}{19:4-12} \\
    \midrule
    耶穌被彼拉多定案 & 耶路撒冷  &       &       & \multicolumn{1}{r|}{23:24} & \multicolumn{1}{p{4.055em}|}{19:13-15} \\
    \midrule
    耶穌被交給人釘死 & 耶路撒冷  & \multicolumn{1}{p{4.055em}|}{27:26下} & \multicolumn{1}{p{4.055em}|}{15:15下} & \multicolumn{1}{r|}{23:25} & 19:16 \\
    \midrule
    古利奈人西門代背十字架 & 耶路撒冷  & 27:32:00 & 15:21 & \multicolumn{1}{r|}{23:26} & 19:17 \\
    \midrule
    耶穌勸戒衆婦人 & 耶路撒冷  &       &       & 23:27-31 &  \\
    \midrule
    耶穌被釘十字架 & 各各他   & \multicolumn{1}{p{4.055em}|}{27:33-35上} & \multicolumn{1}{p{4.055em}|}{3:22，24上，25} & 23:33上，35-38 & 19:18 \\
    \midrule
    耶穌被衆人譏笑 & 各各他   & \multicolumn{1}{p{4.055em}|}{27:37-44} & \multicolumn{1}{p{4.055em}|}{15:26-32} & 23:33下 & \multicolumn{1}{p{4.055em}|}{19:19-22} \\
    \midrule
    兵丁分衣  & 各各他   & \multicolumn{1}{p{4.055em}|}{27:35下-36} & \multicolumn{1}{p{4.055em}|}{15:24下} & 23:34下 & \multicolumn{1}{p{4.055em}|}{19:23-24} \\
    \midrule
    強盜得救  & 各各他   &       &       & 23:39-43 &  \\
    \bottomrule
    \end{tabular}%
     \caption{耶穌被交給人釘死}
  \label{tab:addlabel}%
\end{table}
\subsubsection{彼拉多判耶穌被釘死}
於是，彼拉多把他們所求的那作亂殺人下在監裡的釋放了，把耶穌鞭打了，交給他們去釘十字架。任憑他們的意思行。21 他們就把耶穌帶了去。耶穌背著自己的十字架出來，到了一個地方，名叫髑髏地，希伯來話叫各各他。



\subsubsection{古利奈人背主的十字架-太27:32-33/可15:21-22/路23:26/約19:17}
\paragraph{太27:32-33}他們出來的時候，遇見一個古利奈人，名叫西門，就勉強他同去，好背著耶穌的十字架。到了一個地方名叫各各他，意思就是髑髏地。
\paragraph{可15:21-22}有一個古利奈人西門，就是亞力山大和魯孚的父親，從鄉下來，經過那地方，他們就勉強他同去，好背著耶穌的十字架。
他們帶耶穌到了各各他地方（各各他翻出來就是髑髏地），
\paragraph{路23:26}帶耶穌去的時候，有一個古利奈人西門，從鄉下來；他們就抓住他，把十字架擱在他身上，叫他背著跟隨耶穌。
\paragraph{約19:17}
他們就把耶穌帶了去。耶穌背著自己的十字架出來，到了一個地方，名叫髑髏地，希伯來話叫各各他。
\paragraph{综合}
帶耶穌去的時候，有一個古利奈人西門，就是亞歷山大和魯孚的父親，從鄉下來，經過那地方。他們就勉強他同去，好背著耶穌的十字架。他們就抓住他，把十字架擱在他身上，叫他背著跟隨耶穌。

\subsubsection{耶穌勸眾婦女不要哀哭-路23:27-31}

有許多百姓跟隨耶穌，內中有好些婦女；婦女們為他號咷痛哭。
耶穌轉身對他們說：耶路撒冷的女子，不要為我哭，當為自己和自己的兒女哭。
因為日子要到，人必說：不生育的，和未曾懷胎的，未曾乳養嬰孩的，有福了！
那時，人要向大山說：倒在我們身上！向小山說：遮蓋我們！
這些事既行在有汁水的樹上，那枯乾的樹將來怎麼樣呢？

\subsubsection{耶穌的罪狀-太27:37/可15:25-26/路23:38/約19:19-22}
\paragraph{太27:37}
在他頭以上安一個牌子，寫著他的罪狀，說：這是猶太人的王耶穌。
\paragraph{可15:25-26}
釘他在十字架上是巳初的時候。
在上面有他的罪狀，寫的是：猶太人的王。
\paragraph{路23:38}
在耶穌以上有一個牌子（有古卷在此有：用希利尼、羅馬、希伯來的文字）寫著：這是猶太人的王。
\paragraph{約19:19-22}
彼拉多又用牌子寫了一個名號，安在十字架上，寫的是：猶太人的王，拿撒勒人耶穌。有許多猶太人念這名號；因為耶穌被釘十字架的地方與城相近，並且是用希伯來、羅馬、希利尼三樣文字寫的。猶太人的祭司長就對彼拉多說：不要寫猶太人的王，要寫他自己說：我是猶太人的王。彼拉多說：我所寫的，我已經寫上了。

\subsubsection{兵丁拿苦膽調和的酒給耶穌喝-太27:34/可15:23}
\paragraph{太27：34}
兵丁拿苦膽調和的酒給耶穌喝。他嘗了，就不肯喝。
\paragraph{可15：23}
拿沒藥調和的酒給耶穌，他卻不受。

\subsubsection{兵丁分耶穌的衣服-太27:35-36/可15:24/路23:34a/約19:23-24}
\paragraph{太27:35-36}
他們既將他釘在十字架上，就拈鬮分他的衣服，又坐在那裡看守他。
\paragraph{可15:24}
於是將他釘在十字架上，拈鬮分他的衣服，看是誰得什麼。
\paragraph{路23:34a}
兵丁就拈鬮分他的衣服。
\paragraph{約19:23-24}
兵丁既然將耶穌釘在十字架上，就拿他的衣服分為四分，每兵一分；又拿他的裡衣，這件裡衣原來沒有縫兒，是上下一片織成的。
他們就彼此說：我們不要撕開，只要拈鬮，看誰得著。這要應驗經上的話說：他們分了我的外衣，為我的裡衣拈鬮。兵丁果然做了這事。拈鬮分他的衣服。看是誰得什麼。又坐在那裡看守他。

\subsubsection{同釘的兩個強盜-太27:38,44/可15:27-28,32b/路23:32-33，39-43/約19:18}
\paragraph{太27:38,44}
當時，有兩個強盜和他同釘十字架，一個在右邊，一個在左邊。
那和他同釘的強盜也是這樣的譏誚他。
\paragraph{可15:27-28,32b}
他們又把兩個強盜和他同釘十字架，一個在右邊，一個在左邊。（有古卷在此有：這就應了經上的話說：他被列在罪犯之中。）

那和他同釘的人也是譏誚他。
\paragraph{路23:32-33，39-43}
又有兩個犯人，和耶穌一同帶來處死。
到了一個地方，名叫髑髏地，就在那裡把耶穌釘在十字架上，又釘了兩個犯人：一個在左邊，一個在右邊。

那同釘的兩個犯人有一個譏誚他，說：你不是基督嗎？可以救自己和我們吧！
那一個就應聲責備他，說：你既是一樣受刑的，還不怕神嗎？
我們是應該的，因我們所受的與我們所做的相稱，但這個人沒有做過一件不好的事。
就說：耶穌啊，你得國降臨的時候，求你記念我！
耶穌對他說：我實在告訴你，今日你要同我在樂園裡了。
\paragraph{約19:18}
他們就在那裡釘他在十字架上，還有兩個人和他一同釘著，一邊一個，耶穌在中間。

\subsubsection{耶穌爲人代禱}
當下耶穌說：「父啊，赦免他們！因為他們所做的他們不曉得。」

\subsubsection{衆人的反應-太27:39-43/可15:29-32a/路23:34-37}
\begin{table}[htbp]
  \centering 
    \begin{tabular}{|p{4.225em}|p{35.955em}|}
    \toprule
    從那裡經過的人 & 辱罵他，搖著頭說：「咳！你這拆毀聖殿、三日又建造起來的，可以救自己，你如果是神的兒子，從十字架上下來吧！」 \\
    \midrule
    祭司長和文士 & 戲弄他，彼此說：「他救了別人，不能救自己。以色列的王基督，現在可以從十字架上下來，叫我們看見就信了！他倚靠神，神若喜悅他，現在可以救他！因為他曾說：『我是神的兒子。』」 \\
    \midrule
    強盜    & 那和他同釘的強盜也是譏誚他。 \\
    \midrule
    百姓    & 百姓站在那裡觀看。 \\
    \midrule
    官府    & 嗤笑他「他救了別人，他若是基督，神所揀選的，可以救自己吧！」 \\
    \midrule
    兵丁    & 戲弄他，上前拿醋送給他喝「你若是猶太人的王，可以救自己吧！」 \\
    \bottomrule
    \end{tabular}%
     \caption{耶穌被衆人譏笑}
  \label{tab:addlabel}%
\end{table}

\paragraph{太27:39-43}
從那裡經過的人譏誚他，搖著頭，說：
你這拆毀聖殿、三日又建造起來的，可以救自己吧！你如果是神的兒子，就從十字架上下來吧！
祭司長和文士並長老也是這樣戲弄他，說：
他救了別人，不能救自己。他是以色列的王，現在可以從十字架上下來，我們就信他。
他倚靠神，神若喜悅他，現在可以救他；因為他曾說：我是神的兒子。
\paragraph{可15:29-32a}
從那裡經過的人辱罵他，搖著頭說：咳！你這拆毀聖殿、三日又建造起來的，
可以救自己，從十字架上下來吧！
祭司長和文士也是這樣戲弄他，彼此說：他救了別人，不能救自己。
以色列的王基督，現在可以從十字架上下來，叫我們看見，就信了。
\paragraph{路23:34-37}
當下耶穌說：父啊！赦免他們；因為他們所做的，他們不曉得。
百姓站在那裡觀看。官府也嗤笑他，說：他救了別人；他若是基督，神所揀選的，可以救自己吧！
兵丁也戲弄他，上前拿醋送給他喝，
說：你若是猶太人的王，可以救自己吧！
\paragraph{综合}
29 從那裡經過的人辱罵他，搖著頭說：「咳！你這拆毀聖殿、三日又建造起來的， 30 可以救自己，你如果是神的兒子，從十字架上下來吧！」 31 祭司長和文士也是這樣戲弄他，彼此說：「他救了別人，不能救自己。 32 以色列的王基督，現在可以從十字架上下來，叫我們看見就信了！他倚靠神，神若喜悅他，現在可以救他！因為他曾說：『我是神的兒子。』」那和他同釘的強盜也是譏誚他。百姓站在那裡觀看。官府也嗤笑他，說：「他救了別人，他若是基督，神所揀選的，可以救自己吧！」 36 兵丁也戲弄他，上前拿醋送給他喝， 37 說：「你若是猶太人的王，可以救自己吧！」

\subsubsection{耶穌托付門徒照看自己的母親-約19:25-27}
%\paragraph{約19：17-37}
站在耶穌十字架旁邊的，有他母親與他母親的姊妹，並革羅罷的妻子馬利亞，和抹大拉的馬利亞。耶穌見母親和他所愛的那門徒站在旁邊，就對他母親說：母親（原文作婦人），看，你的兒子！又對那門徒說：看，你的母親！從此，那門徒就接他到自己家裡去了。

\subsubsection{耶穌渴了-太27:45-50/可15:33-36/約19:28-30}
\paragraph{太27:45-50}
從午正到申初，遍地都黑暗了。
約在申初，耶穌大聲喊著說：以利！以利！拉馬撒巴各大尼？就是說：我的神！我的神！為什麼離棄我？
站在那裡的人，有的聽見就說：這個人呼叫以利亞呢！
內中有一個人趕緊跑去，拿海絨蘸滿了醋，綁在葦子上，送給他喝。
其餘的人說：且等著，看以利亞來救他不來。耶穌又大聲喊叫，氣就斷了。
\paragraph{可15:33-36}
從午正到申初，遍地都黑暗了。
申初的時候，耶穌大聲喊著說：以羅伊！以羅伊！拉馬撒巴各大尼？翻出來就是：我的神！我的神！為什麼離棄我？
旁邊站著的人，有的聽見就說：看哪，他叫以利亞呢！
有一個人跑去，把海絨蘸滿了醋，綁在葦子上，送給他喝，說：且等著，看以利亞來不來把他取下。
\paragraph{約19:28-30}
這事以後，耶穌知道各樣的事已經成了，為要使經上的話應驗，就說：我渴了。
有一個器皿盛滿了醋，放在那裡；他們就拿海絨蘸滿了醋，綁在牛膝草上，送到他口。
耶穌嘗（原文作受）了那醋，就說：成了！便低下頭，將靈魂交付神了。

\subsubsection{十架七言}
\begin{table}[htbp]
  \centering
    \begin{tabular}{|p{8.07em}|p{20.645em}|p{7.85em}|}
    \toprule
    對象    & 耶穌的話  & 經文 \\
    \midrule
    釘他的人  & 父啊！赦免他們，因為他們所作的，他們不曉得。 & 路23：34 \\
    \midrule
    同釘強盜  & 我實在告訴你，今日你要同我在樂園了。 & 路23：43 \\
    \midrule
    母親瑪利亞和約翰 & 「母親，看你的兒子。」「看你的母親。」 & 約19：26-27 \\
    \midrule
    \multirow{1}{*}{父神（午正-申初）} & 我的神！我的神！為甚麼離棄我？ & \multirow{1}{*}{太27：46可15:34} \\   
    \midrule
    兵丁    & 我渴了。（喝醋-可15：36/約19:29） & 約19:28 \\
    \midrule
    自己    & 成了！   & 約19:30 \\
    \midrule
    父神    & 父啊！我將我的靈魂交在你手裏。 & 路23：46 \\
    \bottomrule
    \end{tabular}%
   \caption{十架七言}  
  \label{tab:addlabel}%
\end{table}


\subsubsection{耶穌斷氣-太27:51-56/可15:37-38，40-41/路23:44-49}


\paragraph{太27:51-56}忽然，殿裡的幔子從上到下裂為兩半，地也震動，磐石也崩裂，
墳墓也開了，已睡聖徒的身體多有起來的。
到耶穌復活以後，他們從墳墓裡出來，進了聖城，向許多人顯現。
百夫長和一同看守耶穌的人看見地震並所經歷的事，就極其害怕，說：這真是神的兒子了！
有好些婦女在那裡，遠遠的觀看；他們是從加利利跟隨耶穌來服事他的。
內中有抹大拉的馬利亞，又有雅各和約西的母親馬利亞，並有西庇太兩個兒子的母親。
\paragraph{可15:37-38，40-41}
耶穌大聲喊叫，氣就斷了。殿裡的幔子從上到下裂為兩半。
還有些婦女遠遠的觀看；內中有抹大拉的馬利亞，又有小雅各和約西的母親馬利亞，並有撒羅米，
就是耶穌在加利利的時候，跟隨他、服事他的那些人，還有同耶穌上耶路撒冷的好些婦女在那裡觀看。
\paragraph{路23:44-49}
那時約有午正，遍地都黑暗了，直到申初，
日頭變黑了；殿裡的幔子從當中裂為兩半。
耶穌大聲喊著說：父啊！我將我的靈魂交在你手裡。說了這話，氣就斷了。
聚集觀看的眾人見了這所成的事都捶著胸回去了。
還有一切與耶穌熟識的人，和從加利利跟著他來的婦女們，都遠遠的站著看這些事。
\paragraph{综合}
從午正到申初日頭變黑了，，遍地都黑暗了。 46 約在申初，耶穌大聲喊著說：「以利！以利！拉馬撒巴各大尼？」就是說：「我的神！我的神！為什麼離棄我？」 47 站在那裡的人，有的聽見就說：「這個人呼叫以利亞呢！」 48 內中有一個人趕緊跑去，拿海絨蘸滿了醋，綁在葦子上，送給他喝。 49 其餘的人說：「且等著，看以利亞來救他不來。」 50 耶穌又大聲喊叫，氣就斷了。這事以後，耶穌知道各樣的事已經成了，為要使經上的話應驗，就說：「我渴了。」29 有一個器皿盛滿了醋，放在那裡。他們就拿海絨蘸滿了醋，綁在牛膝草上，送到他口。 30 耶穌嘗 c了那醋，就說：「成了！」耶穌大聲喊著說：「父啊，我將我的靈魂交在你手裡！」說了這話，便低下頭，氣就斷了，將靈魂交付神了。

\subsubsection{百夫長承認耶穌-可15:39/路23:47}
\paragraph{可15:39}
對面站著的百夫長看見耶穌這樣喊叫（有古卷沒有喊叫二字）斷氣，就說：這人真是神的兒子！
\paragraph{路23：47}百夫長看見所成的事，就歸榮耀與神，說：這真是個義人！

\subsubsection{殿幔與墳墓裂開等奇事}
忽然，殿裡的幔子從上到下裂為兩半，地也震動，磐石也崩裂， 52 墳墓也開了，已睡聖徒的身體多有起來的。 53 到耶穌復活以後，他們從墳墓裡出來，進了聖城，向許多人顯現。 54 對面站著的百夫長和一同看守耶穌的人看見地震, 耶穌這樣喊叫斷氣，並所經歷的事，就極其害怕歸榮耀於神，說：「這真是神的兒子了！這真是個義人！」聚集觀看的眾人見了這所成的事，都捶著胸回去了。 55 有好些婦女在那裡遠遠地觀看，她們是從加利利跟隨耶穌來服侍他的。 56 內中有抹大拉的馬利亞，又有雅各和約西的母親馬利亞，並有西庇太兩個兒子的母親。並有撒羅米， 41 就是耶穌在加利利的時候，跟隨他、服侍他的那些人；還有同耶穌上耶路撒冷的好些婦女在那裡觀看。

\subsubsection{耶穌肋旁被紮}
猶太人因這日是預備日，又因那安息日是個大日，就求彼拉多叫人打斷他們的腿，把他們拿去，免得屍首當安息日留在十字架上。 32 於是兵丁來，把頭一個人的腿，並與耶穌同釘第二個人的腿，都打斷了。 33 只是來到耶穌那裡，見他已經死了，就不打斷他的腿。 34 唯有一個兵拿槍扎他的肋旁，隨即有血和水流出來。 35 看見這事的那人就作見證，他的見證也是真的，並且他知道自己所說的是真的，叫你們也可以信。 36 這些事成了，為要應驗經上的話說：「他的骨頭一根也不可折斷。」 37 經上又有一句說：「他們要仰望自己所扎的人。」



\subsection{被埋葬-太27:57-61/可15:42-47/路23:50-55/約19:31-42}
\paragraph{太27:57-61}
到了晚上，有一個財主，名叫約瑟，是亞利馬太來的，他也是耶穌的門徒。
這人去見彼拉多，求耶穌的身體；彼拉多就吩咐給他。
約瑟取了身體，用乾淨細麻布裹好，
安放在自己的新墳墓裡，就是他鑿在磐石裡的。他又把大石頭滾到墓門口，就去了。
有抹大拉的馬利亞和那個馬利亞在那裡，對著墳墓坐著。

\paragraph{可15:42-47}
到了晚上，因為這是預備日，就是安息日的前一日，
有亞利馬太的約瑟前來，他是尊貴的議士，也是等候神國的。他放膽進去見彼拉多，求耶穌的身體；
彼拉多詫異耶穌已經死了，便叫百夫長來，問他耶穌死了久不久。
既從百夫長得知實情，就把耶穌的屍首賜給約瑟。
約瑟買了細麻布，把耶穌取下來，用細麻布裹好，安放在磐石中鑿出來的墳墓裡，又滾過一塊石頭來擋住墓門。
抹大拉的馬利亞和約西的母親馬利亞都看見安放他的地方。

\paragraph{路23:50-55}

有一個人名叫約瑟，是個議士，為人善良公義；
眾人所謀所為，他並沒有附從。他本是猶太、亞利馬太城裡素常盼望神國的人。
這人去見彼拉多，求耶穌的身體，
就取下來，用細麻布裹好，安放在石頭鑿成的墳墓裡；那裡頭從來沒有葬過人。
那日是預備日，安息日也快到了。
那些從加利利和耶穌同來的婦女跟在後面，看見了墳墓和他的身體怎樣安放。

\paragraph{約19:31-42}
猶太人因這日是預備日，又因那安息日是個大日，就求彼拉多叫人打斷他們的腿，把他們拿去，免得屍首當安息日留在十字架上。
於是兵丁來，把頭一個人的腿，並與耶穌同釘第二個人的腿，都打斷了。
只是來到耶穌那裡，見他已經死了，就不打斷他的腿。
惟有一個兵拿槍扎他的肋旁，隨即有血和水流出來。
看見這事的那人就作見證─他的見證也是真的，並且他知道自己所說的是真的─叫你們也可以信。
這些事成了，為要應驗經上的話說：他的骨頭一根也不可折斷。
經上又有一句說：他們要仰望自己所扎的人。
這些事以後，有亞利馬太人約瑟，是耶穌的門徒，只因怕猶太人，就暗暗的作門徒。他來求彼拉多，要把耶穌的身體領去。彼拉多允准，他就把耶穌的身體領去了。又有尼哥底母，就是先前夜裡去見耶穌的，帶著沒藥和沉香約有一百斤前來。他們就照猶太人殯葬的規矩，把耶穌的身體用細麻布加上香料裹好了。在耶穌釘十字架的地方有一個園子，園子裡有一座新墳墓，是從來沒有葬過人的。只因是猶太人的預備日，又因那墳墓近，他們就把耶穌安放在那裡。
\paragraph{综合}
這些事以後，到了晚上，有亞利馬太的財主約瑟，是個尊貴的議士，為人善良公義， 51 眾人所謀所為他並沒有附從；他本是猶太亞利馬太城裡素常盼望神國的人。是耶穌的門徒，只因怕猶太人，就暗暗地做門徒，他來放膽進去求彼拉多，要把耶穌的身體領去。詫異耶穌已經死了，便叫百夫長來，問他耶穌死了久不久。 45 既從百夫長得知實情，就把耶穌的屍首賜給約瑟。他就把耶穌的身體領去了。約瑟買了細麻布， 39 又有尼哥迪慕，就是先前夜裡去見耶穌的，帶著沒藥和沉香約有一百斤前來。 40 他們就照猶太人殯葬的規矩，把耶穌的身體用乾淨細麻布加上香料裹好了。 41 在耶穌釘十字架的地方有一個園子，園子裡有一座新的在磐石中鑿成的墳墓，是從來沒有葬過人的。 42 只因是猶太人的預備日，安息日也快到了, 又因那墳墓近，他們就把耶穌安放在那裡。又滾過一塊石頭來擋住墓門。那些從加利利和耶穌同來的婦女又滾過一塊石頭來擋住墓門。又滾過一塊石頭來擋住墓門。抹大拉的馬利亞和約西的母親馬利亞都跟在後面，看見了墳墓和他的身體怎樣安放。對著墳墓坐著。 56 她們就回去，預備了香料香膏。


